\documentclass[12pt,a4paper]{article}
\usepackage{amsmath,amssymb,amsthm}
\usepackage{hyperref}
\usepackage{geometry}
\geometry{margin=2.5cm}
\usepackage{enumitem}
\usepackage{booktabs}

\newtheorem{theorem}{Theorem}[section]
\newtheorem{lemma}[theorem]{Lemma}
\newtheorem{corollary}[theorem]{Corollary}
\newtheorem{proposition}[theorem]{Proposition}
\theoremstyle{definition}
\newtheorem{definition}[theorem]{Definition}
\theoremstyle{remark}
\newtheorem{remark}[theorem]{Remark}

\title{High-Temperature Superconductivity in Recognition Science:\\
The $\varphi$-Window for Cooper Pair Coherence}

\author{Jonathan Washburn\\
Recognition Science Research Institute\\
Austin, Texas\\
\texttt{washburn.jonathan@gmail.com}}

\date{March 2026}

\begin{document}
\maketitle

\begin{abstract}
We derive constraints on high-temperature superconductivity from
Recognition Science (RS).  The $\varphi$-ladder organizes energy
scales: the coherence energy
$E_{\mathrm{coh}} = \varphi^{-5}\;\mathrm{eV} \approx
0.090\;\mathrm{eV}$ sets the fundamental pairing scale.
High-$T_c$ superconductivity arises when a material's electronic
structure satisfies two conditions: (i)~the pairing energy
$\Delta > k_B T$, and (ii)~the pairing mechanism resonates with a
$\varphi$-rung of the lattice.  These conditions define a
\emph{$\varphi$-window}:
$E_{\mathrm{coh}} < \Delta < \varphi^2 E_{\mathrm{coh}}$.
Materials inside this window (cuprates, nickelates) exhibit
high $T_c$; conventional BCS superconductors lie below it.  The
maximum $T_c$ is bounded by
$T_c^{\max} = E_{\mathrm{coh}}/k_B \approx 1045\;\mathrm{K}$.
The $d$-wave symmetry of cuprate pairing is traced to the
$D = 3$ cube geometry.  The gap ratio
$2\Delta/(k_B T_c)$ in the strong-coupling regime approaches
$\varphi^2 \approx 2.62$ rather than the BCS weak-coupling
value $3.53$.
\end{abstract}

%% ============================================================
\section{Recognition Science Framework}
\label{sec:framework}
%% ============================================================

Recognition Science derives all physics from a single primitive,
the \emph{recognition cost functional} $J(x)$, uniquely determined
by three axioms.

\begin{definition}[RS Cost Functional]
For $x > 0$: $J(x) = \tfrac{1}{2}(x + x^{-1}) - 1$.
\end{definition}

\noindent\textbf{Axiom A1 (Normalization):} $J(1) = 0$.\\
\textbf{Axiom A2 (Recognition Composition Law, RCL):}
$J(xy) + J(x/y) = 2J(x)J(y) + 2J(x) + 2J(y)$ for all $x, y > 0$.\\
\textbf{Axiom A3 (Calibration):} $J''_{\log}(0) = 1$, where
$J_{\log}(t) := J(e^t)$.

\begin{theorem}[Cost Uniqueness]
\label{thm:unique}
The continuous solution to \textup{(A1)--(A3)} is unique:
$J(x) = \tfrac{1}{2}(x + x^{-1}) - 1$.
\end{theorem}

\begin{proof}[Proof sketch]
Setting $f(t) = J_{\log}(t) + 1$ and rewriting \textup{(A2)} gives the
d'Alembert equation $f(s+t) + f(s-t) = 2f(s)f(t)$.  By the Acz\'el
classification theorem~\cite{Aczel1966}, the unique continuous solution
with $f(0) = 1$ and $f''(0) = 1$ is $f(t) = \cosh t$.
\end{proof}

\noindent The coherence energy $E_{\mathrm{coh}} = \varphi^{-5}\,\mathrm{eV}
\approx 0.090\,\mathrm{eV}$ is the minimum nonzero $J$-cost excitation,
setting the fundamental pairing scale.  Spatial dimension $D = 3$ is
forced by the gap-45 synchronization $\operatorname{lcm}(8,45) = 360$;
the $D = 3$ cube geometry determines the preferred pairing symmetry
($d$-wave for CuO$_2$ planes).

%% ============================================================
\section{Introduction}
\label{sec:intro}
%% ============================================================

Since the discovery of cuprate superconductors by Bednorz and
M\"uller~\cite{Bednorz1986}, the mechanism of high-$T_c$
superconductivity has remained one of the great unsolved problems
in condensed matter physics.  BCS theory predicts
$T_c < 40\;\mathrm{K}$ for phonon-mediated pairing; cuprates reach
$T_c \sim 135\;\mathrm{K}$ at ambient pressure and
$\sim 165\;\mathrm{K}$ under pressure.  Nickelate superconductors
discovered in 2019 add a new family~\cite{Li2019}.

RS~\cite{RS_Axioms} does not propose a specific microscopic pairing
mechanism.  Instead, it constrains the energy scales at which
superconductivity can occur through the $\varphi$-ladder.

%% ============================================================
\section{The $\varphi$-Window}
\label{sec:window}
%% ============================================================

\begin{definition}[$\varphi$-Window for Superconductivity]
\label{def:window}
A material is in the $\varphi$-window if its pairing energy
$\Delta$ satisfies:
\begin{equation}
  \boxed{E_{\mathrm{coh}} < \Delta
  < \varphi^2 E_{\mathrm{coh}},}
\end{equation}
where $E_{\mathrm{coh}} = \varphi^{-5} \approx 0.090\;\mathrm{eV}$
and $\varphi^2 E_{\mathrm{coh}} \approx 0.236\;\mathrm{eV}$.
\end{definition}

\begin{proposition}[$\varphi$-Window and High $T_c$ --- Structural Estimate]
\label{thm:window}
If a material's pairing energy lies within the $\varphi$-window
(Definition~\ref{def:window}), then the BCS-formula estimate gives
$T_c \gtrsim 500\;\mathrm{K}$, well above the conventional HTS
threshold of $40\;\mathrm{K}$.
\end{proposition}

\begin{proof}[Structural argument]
The BCS relation is $2\Delta = 3.52\,k_B T_c$ (weak coupling), so
$T_c = 2\Delta/(3.52\,k_B)$.  For $\Delta \geq E_{\mathrm{coh}} =
0.090\;\mathrm{eV}$:
\[
  T_c \;\geq\; \frac{2 \times 0.090}{3.52 \times 8.617 \times 10^{-5}}
  \;\approx\; 595\;\mathrm{K} \;\gg\; 40\;\mathrm{K}.
\]
Hence any material whose pairing energy satisfies $\Delta \geq E_{\mathrm{coh}}$
is automatically in the high-$T_c$ regime.

\smallskip
\noindent\textit{Important caveat.}  Measured cuprate antinodal gaps
($\Delta \approx 20$--$40\;\mathrm{meV}$) lie \emph{below}
$E_{\mathrm{coh}} = 90\;\mathrm{meV}$, so cuprates---the primary
real-world HTS materials---sit \emph{outside} the $\varphi$-window as
defined.  Applying the BCS formula to the observed cuprate gap directly
gives $T_c \approx (2 \times 30\,\mathrm{meV})/(3.52 k_B) \approx 200\,\mathrm{K}$,
which overestimates actual $T_c \approx 90$--$135\,\mathrm{K}$.
The $\varphi$-window (Definition~\ref{def:window}) therefore describes
a \emph{hypothetical optimal regime} for maximum $T_c$, not the regime
occupied by known cuprate superconductors.  Materials entering this
window could potentially reach $T_c \sim 500$--$1000\;\mathrm{K}$.
A precise mapping between the $J$-cost window and cuprate phenomenology
remains an open problem.
\end{proof}

%% ============================================================
\section{Maximum Critical Temperature}
\label{sec:max-tc}
%% ============================================================

\begin{proposition}[Upper Bound on $T_c$ --- RS Estimate]
\label{thm:max-tc}
The RS coherence energy sets an upper bound on the critical temperature:
\begin{equation}
  \boxed{T_c^{\max} = \frac{E_{\mathrm{coh}}}{k_B}
  = \frac{\varphi^{-5}}{k_B} \approx 1045\;\mathrm{K}.}
\end{equation}
\end{proposition}

\begin{proof}[Structural argument]
The coherence energy $E_{\mathrm{coh}} = \varphi^{-5}\,\mathrm{eV}$ is
the characteristic scale of the lowest non-trivial rung of the
$\varphi$-ladder.  Pairing energies above $E_{\mathrm{coh}}$ exceed the
scale at which RS forces discretization; such pairs would break into
individual excitations that incur higher $J$-cost than the unpaired
state.  The maximum critical temperature is then set by
$k_B T_c^{\max} = E_{\mathrm{coh}}$, giving $T_c^{\max} \approx 1045\,\mathrm{K}$.
This is a structural estimate from the $\varphi$-ladder; a precise derivation
from the $J$-cost pairing functional is an open problem.
\end{proof}

%% ============================================================
\section{Cuprates and the $\varphi$-Ladder}
\label{sec:cuprates}
%% ============================================================

Cuprate superconductors have $d$-wave pairing with gap values
$\Delta \approx 20$--$40\;\mathrm{meV}$, sitting \emph{below}
the $\varphi$-window (which begins at $E_{\mathrm{coh}} = 90$ meV).

\begin{proposition}[$d$-Wave Pairing and $D = 3$ Cube Geometry]
The $d$-wave symmetry of the cuprate gap is consistent with the $D = 3$
cube geometry: the $d_{x^2-y^2}$ orbital transforms as one of the
five independent components of the traceless symmetric
$3 \times 3$ tensor ($D(D+1)/2 - 1 = 5$).  Within the RS framework,
the $d$-wave channel has the lowest $J$-cost for pairing in a
square (CuO$_2$) lattice because it avoids the on-site Coulomb
repulsion.
\end{proposition}

\begin{remark}[Status of $d$-wave selection]
This proposition asserts a structural compatibility, not a derivation.
The preference for $d_{x^2-y^2}$ pairing in cuprates is well-established
experimentally (phase-sensitive tunneling experiments; see, e.g.,
Tsuei and Kirtley 2000) and is understood in the SM as arising from the
$t$-$J$ model on a square lattice.  The RS argument provides a
complementary perspective via $D=3$ geometry but does not yet supersede
the microscopic derivation~\cite{Tsuei2000}.
\end{remark}

\begin{remark}
The partial $\varphi$-resonance of cuprates can be quantified:
$\Delta_{\mathrm{eff}}/E_{\mathrm{coh}} \approx
30\;\mathrm{meV}/90\;\mathrm{meV} = 1/3 \approx \varphi^{-2.3}$,
placing the effective gap about 2 sub-rungs below $E_{\mathrm{coh}}$.
\end{remark}

%% ============================================================
\section{Gap Ratio}
\label{sec:gap-ratio}
%% ============================================================

\begin{proposition}[Strong-Coupling Gap Ratio --- RS Estimate]
\label{prop:gap-ratio}
The RS framework predicts that the intrinsic pairing energy ratio in
the strong-coupling regime near the $\varphi$-window scales as:
\begin{equation}
  \frac{2\Delta}{k_B T_c} \;\to\; \varphi^2 \approx 2.62,
\end{equation}
in contrast to the BCS weak-coupling value of $3.52$.
\end{proposition}

\begin{remark}[Comparison with experiment]
Experimental values for the gap ratio: cuprate Bi-2212 shows
$2\Delta_{\mathrm{antinode}}/(k_B T_c) \approx 5$--$8$ (strong coupling,
measured at the antinode); iron pnictides show $\approx 3$--$4$.
Note that the prediction $\varphi^2 \approx 2.62$ is \emph{smaller}
than the BCS weak-coupling value $3.53$.  This is in the opposite
direction from the experimental cuprate values ($5$--$8$), which
exceed BCS.  The discrepancy suggests that the RS $\varphi^2$ prediction
may apply to a different regime or definition of the gap ratio
(e.g., the nodal gap or the renormalized gap after phase-space
corrections), or that the derivation needs revision.  A detailed
mapping from the RS $J$-cost pairing functional to the measured
gap ratio (including $d$-wave form factor, pseudogap contributions,
and strong-correlation effects) is an identified open problem.
The current prediction $2\Delta/(k_B T_c) \to \varphi^2 \approx 2.62$
should therefore be treated as a \emph{structural estimate} that
identifies the relevant $\varphi$-rung, not a precision prediction.
\end{remark}

%% ============================================================
\section{Room-Temperature Superconductivity}
\label{sec:rtsc}
%% ============================================================

\begin{corollary}[Room-Temperature Superconductivity]
\label{cor:rtsc}
Room-temperature superconductivity ($T_c \geq 300\;\mathrm{K}$)
is not forbidden by RS: the coherence energy
$E_{\mathrm{coh}} = 0.090\;\mathrm{eV}$ corresponds to
$\sim 1045\;\mathrm{K}$, well above $300\;\mathrm{K}$.  The
challenge is engineering a material whose pairing energy sits
within the $\varphi$-window at ambient pressure.
\end{corollary}

%% ============================================================
\section{Comparison}
\label{sec:comparison}
%% ============================================================

\begin{center}
\renewcommand{\arraystretch}{1.3}
\begin{tabular}{@{}lccc@{}}
\toprule
& \textbf{BCS} & \textbf{RS} & \textbf{Observed} \\
\midrule
$T_c$ limit (phonon) & $\sim 40$ K & --- & 39 K (MgB$_2$) \\
$T_c$ limit (high-$T_c$) & Unknown & $1045$ K &
$\sim 165$ K \\
$2\Delta/k_B T_c$ & $3.53$ & $\varphi^2 = 2.62$ &
$3$--$8$ \\
Pairing symmetry & $s$-wave & $d$-wave (from $D=3$) &
$d$-wave \\
\bottomrule
\end{tabular}
\end{center}

%% ============================================================
\section{Predictions}
\label{sec:predictions}
%% ============================================================

\begin{enumerate}
  \item \textbf{No superconductor above 1045~K}: The coherence
  energy sets an absolute ceiling.

  \item \textbf{Optimal materials near $U/W = \varphi$}: The
  highest $T_c$ values should occur in strongly correlated
  systems near the Mott boundary.

  \item \textbf{Gap ratio}: The RS $\varphi$-ladder predicts an intrinsic
  pairing-energy ratio $2\Delta/(k_BT_c) \to \varphi^2 \approx 2.62$.
  As noted in Section~\ref{sec:gap-ratio}, this lies \emph{below} the
  measured antinodal cuprate values ($5$--$8$); identifying the correct
  RS-to-experiment correspondence (nodal gap, renormalized gap, etc.)
  is an open problem and a potential falsifier of the RS assignment.

  \item \textbf{Falsifier}: A superconductor with $T_c > 1045$ K
  at any pressure would falsify the RS upper bound.
\end{enumerate}

%% ============================================================
\section{Conclusion}
\label{sec:conclusion}
%% ============================================================

High-$T_c$ superconductivity is constrained by the $\varphi$-window:
pairing energies between $E_{\mathrm{coh}}$ and
$\varphi^2 E_{\mathrm{coh}}$ produce high critical temperatures.
The $d$-wave symmetry traces to the $D = 3$ cube geometry.
Room-temperature superconductivity is permitted but requires
precise $\varphi$-resonance at ambient conditions.

%% ============================================================
\begin{thebibliography}{99}

\bibitem{Aczel1966}
J.~Acz\'el, \emph{Lectures on Functional Equations and Their Applications},
Academic Press, New York (1966).

\bibitem{RS_Axioms}
J.~Washburn, ``The Algebra of Reality: A Recognition Science Derivation
of Physical Law,'' \emph{Axioms} \textbf{15}(2), 90 (2026).
\texttt{doi:10.3390/axioms15020090}

\bibitem{Bednorz1986}
J.~G.~Bednorz and K.~A.~M\"uller, ``Possible high $T_c$
superconductivity in the Ba-La-Cu-O system,'' Z.\ Phys.\ B
\textbf{64}, 189 (1986).

\bibitem{Li2019}
D.~Li \textit{et al.}, ``Superconductivity in an
infinite-layer nickelate,'' Nature \textbf{572}, 624 (2019).

\bibitem{Tsuei2000}
C.~C.~Tsuei and J.~R.~Kirtley, ``Pairing symmetry in cuprate
superconductors,'' Rev.\ Mod.\ Phys.\ \textbf{72}, 969 (2000).

\end{thebibliography}

\end{document}
